\section{Background}
\label{sec:greybox-background}
In this section, we introduce some necessary relevant background and prior works in this field.

% \subsection{\zksnark}

% Zero-knowledge proofs allow a prover to convince a verifier that a statement is true without revealing any information beyond the validity of the statement \cite{groth2016size}. Among different zero-knowledge proof constructions, zero-knowledge succinct non-interactive arguments of knowledge (zk-SNARKs) have gained significant attention because of their efficiency and practicality. zk-SNARKs provide three important properties. First, zero-knowledge ensures that the proof does not reveal the prover's private witness. Second, succinctness enables compact proofs and efficient verification. Third, non-interactivity allows the prover to send a single proof to the verifier without requiring multiple rounds of communication. These properties make zk-SNARKs suitable for applications that require both privacy and efficiency, such as anonymous transactions, verifiable computation, and privacy-preserving identity systems. Their adoption in blockchain systems, including Zcash and Ethereum, further demonstrates their practical importance \cite{zcash_protocol_specification, fan2025byzsfl, kosba2016hawk, bowe2020zexe}. A Rank-1 Constraint System (R1CS) \cite{ben2019aurora} represents a computation as a set of quadratic constraints over finite-field variables and serves as a common intermediate representation for expressing the statements and witnesses used by many zkSNARK constructions.

%\clg{We only handle specific classes/types of SNARKs right?  Should probably make this clear that this specific work at least doesn't handle ALL SNARKs out there.}

\subsection{Fuzzing}

Fuzzing is an automated software testing technique that generates random or mutated inputs and feeds them to a program to uncover crashes, memory errors, and security vulnerabilities. It has been widely used to detect bugs in many types of software, including cryptographic programs, by exploring unexpected edge cases and triggering abnormal behavior \cite{aflplusplus_afl_fuzz_approach, google_fuzzing_docs}. Differential fuzzing is a specialized fuzzing technique that detects logic errors by comparing the behavior of multiple implementations of the same specification. Unlike traditional fuzzing, which primarily detects memory corruption and crashes, differential fuzzing focuses on output inconsistencies when different programs process the same inputs. Several fuzzing tools have been developed for security testing, as discussed in \autoref{subsec:greybox-prior_work}. Fuzzing techniques can be broadly categorized into white-box, black-box, and grey-box fuzzing.

\parhead{White-box fuzzing} has access to the program's source code and internal structure. It can use techniques such as symbolic execution and constraint solving to systematically explore execution paths \cite{godefroid2007random}.

\parhead{Black-box fuzzing} tests a program without access to its source code or internal structure. It relies only on external inputs and outputs to detect vulnerabilities. Black-box fuzzing generates random or mutated inputs and observes program behavior to identify crashes or unexpected results, but it lacks detailed execution-path awareness \cite{godefroid2007random}.

\parhead{Grey-box fuzzing} operates with partial knowledge of the program's internal behavior without requiring source code. It often uses lightweight instrumentation or feedback mechanisms to guide input generation and improve path exploration \cite{wang2020sok}.

\subsection{Differential Fuzzing}
Differential fuzzing evaluates a target program by repeatedly executing it alongside a test-vector program that implements the same specification. Both programs receive the same generated inputs, and their outputs are compared to identify inconsistencies that may indicate implementation errors or security vulnerabilities.

For cryptographic programs, providing correct test vectors presents unique challenges due to the structured and deterministic nature of cryptographic algorithms. Some cryptographic algorithms such as RSA and SHA256 have standardized test vector program, which can be used in differential fuzzing.  Known test vectors derived from standard cryptographic test suites, such as NIST's Cryptographic Algorithm Validation Program (CAVP), serve as valuable seed inputs \cite{nist_cryptographic_algorithm_validation_program}. Differential fuzzing further enhances this process by comparing outputs from multiple cryptographic implementations, such as OpenSSL, BoringSSL, and LibreSSL, to detect inconsistencies that may indicate implementation flaws. However, a complex cryptographic protocol such as \zkp does not have an existing test vector program. In order to develop a differential fuzzing approach to detect \logicerror, the differential fuzzing needs to obtain a test vector for \zkp program.


\subsection{Taint Tracking}

Taint tracking is a program analysis technique that monitors how sensitive or untrusted data flows through a program. It is commonly used to detect security vulnerabilities such as information leaks, injection attacks, and unsafe memory operations. Taint tracking marks selected inputs as tainted and propagates taint as those inputs influence registers, memory locations, variables, or program outputs. This allows an analysis tool to determine whether untrusted data reaches security-critical operations, such as system calls, memory accesses, or cryptographic computations. Taint tracking can be performed dynamically, where taint propagation is analyzed during program execution, or statically, where the program is analyzed without execution \cite{yan2017sok}.


\subsection{Prior Works}
\label{subsec:greybox-prior_work}

\begin{table}[ht]
\caption{Comparison of \greyboxfuzz and a selection of cryptographic and generic fuzzing tools}
\label{tab:fuzzer_comparison}
\vspace*{6pt}
\centering
% \scriptsize{
\begin{tabular}{c||c|c|c|c|c}
\hline
Fuzzer
& \rotatebox{70}{AFL \cite{google_afl}}
& \rotatebox{70}{CDF \cite{aumasson2017automated}}
& \rotatebox{70}{SNARKProbe \cite{fan2024snarkprobe}}
& \rotatebox{70}{ZKFUZZ \cite{takahashi2026zkfuzz}}
& \rotatebox{70}{\greyboxfuzz} \\ \hline\hline
Blackbox Binary & \CIRCLE & \CIRCLE & \Circle & \Circle & \CIRCLE \\ \hline
Support \zksnark & \CIRCLE & \Circle & \CIRCLE & \CIRCLE & \CIRCLE \\ \hline
Software Error & \CIRCLE & \CIRCLE & \CIRCLE & \Circle & \CIRCLE \\ \hline
Crypto Error & \Circle & \CIRCLE & \CIRCLE & \RIGHTcircle & \CIRCLE \\ \hline
Error Locator & \RIGHTcircle & \Circle & \Circle & \Circle & \CIRCLE \\ \hline
\end{tabular}
% }
\end{table}


Several fuzzing tools can be used to evaluate the security of cryptographic programs, and some of them are specifically designed for cryptographic software to improve testing performance and detection capability. However, these techniques still have important limitations. Some tools do not support \zksnark programs or cannot detect \logicerror because they mainly focus on crashes, memory errors, or other generic software failures. Other fuzzing tools that support cryptographic protocols often require access to source code, which reduces their applicability to real-world binary-only programs. As shown in \autoref{tab:fuzzer_comparison}, existing approaches also provide limited support for locating the source of detected errors.

\begin{table}[ht]
\caption{Comparison of the focus of \zksnark security tools}
\label{tab:snark_comparison}
\vspace*{6pt}
\centering
\begin{tabular}{c||c|c|c}
\hline
Tool & Circuit Error & Frontend Error & Backend Error \\ \hline\hline
Circomspect \cite{trailofbits2023circomspect} & \CIRCLE & \Circle & \Circle \\ \hline
ZKAP \cite{wen2024practical} & \CIRCLE & \Circle & \Circle \\ \hline
Korrekt \cite{soureshjani2023automated} & \CIRCLE & \Circle & \Circle \\ \hline
Coda \cite{liu2024certifying} & \CIRCLE & \Circle & \Circle \\ \hline
Ecne \cite{ecne} & \CIRCLE & \Circle & \Circle \\ \hline
Picus \cite{pailoor2023automated} & \CIRCLE & \Circle & \Circle \\ \hline
Aleo \cite{coglio2023compositional} & \CIRCLE & \Circle & \Circle \\ \hline
OWBB23 \cite{ozdemir2025bounded} & \Circle & \CIRCLE & \Circle \\ \hline
SNARKProbe \cite{fan2024snarkprobe} & \CIRCLE & \Circle & \CIRCLE \\ \hline
CIVER \cite{isabel2024scalable} & \CIRCLE & \Circle & \Circle \\ \hline
ZKFUZZ \cite{takahashi2026zkfuzz} & \CIRCLE & \Circle & \Circle \\ \hline
\end{tabular}
\end{table}

A growing body of work has also studied \zksnark security analysis \cite{chaliasos2024sok}, especially circuit-level checking. These tools mainly focus on circuit errors, including missing constraints, under-constrained signals, incorrect range checks, and unsafe gadget usage. Although circuit-level analysis is important, it does not cover the full deployed workflow, where errors may also arise in frontend input conversion, setup, proof generation, proof verification, or the interaction between components. As summarized in \autoref{tab:snark_comparison}, prior \zksnark security tools mainly focus on circuit-level errors and provide limited support for frontend or backend implementation errors.

As a result, existing tools are insufficient for evaluating deployed \zksnark applications in binary-only settings. They either focus on generic software failures, require source-code or circuit-level access, or analyze only part of the \zksnark workflow. \greyboxfuzz addresses this gap by testing binary prover and verifier programs end to end against the intended cryptographic statement, while also providing structured input generation, cryptographic coverage monitoring, and component-level error localization.